\section{Related work}\label{sec:related}
\paragraph{Silent tool failures.} Benchmarks inject tool perturbations and find implicit failures
harder than explicit ones~\cite{zhu2026tools}, and agents often fail to detect results that come back
well formed but wrong~\cite{kraishan2026loud}, or, under repeated poisoning, adopt them
even after checking~\cite{anon2026tools}. AgentFaultBench checks recovery against the resulting state as well as the
agent's response~\cite{gupta2026agentfaultbench}; RITE studies silent wrong values in multi-step
e-commerce agents~\cite{zhang2026rite}. A verification-first monitor that issues corrective directives
improves recovery from tool failures on a real filesystem~\cite{kaynar2026tools}, and recoveries after
partial executions can themselves harm the external state~\cite{trofimov2026agent}. We share their
concern with faults that look like success, and add a comparison of triggers with recovery held fixed.

\paragraph{Read-back and contracts.} Runtime monitors of action preconditions and postconditions drive
replanning in classical planning~\cite{cardoso2026rvreplan}. For LLM agents, verified tool calls address
non-atomic behaviour such as delayed visibility and partial updates~\cite{mansoor2026verified}, and
LIMBO asks whether exactly-once behaviour belongs in the model, the harness or the tool
contract~\cite{li2026where}. Outcome Monitors check results against outcome contracts and issue a
receipt on a violation~\cite{panthi2026outcome}; SilentGuard enforces invariants in the
harness~\cite{dip2026invariantaware}; RECEIPT measures and repairs the assumption that a call returns
with its action finished~\cite{yang2026receipt}. Our read-back and contract triggers are simplified
versions of these ideas, and our read-back is idealised (Section~\ref{sec:method}), so that the
comparison is about the trigger and not its engineering.

\paragraph{When to intervene.} Recovery can rescue a failing trajectory or disrupt a good
one~\cite{xiao2026harnesses}, and agents often recover on their own if left alone~\cite{duan2026intervene}.
Uncertainty is a common but unproven signal for when to correct an agent~\cite{li2026uncertainty}, a
calibrated failure score need not say whether intervening helps~\cite{zhang2026calibration}, and a gate
that only replaces a proposal can act only where the base policy was already
wrong~\cite{sethi2026consistency}. Middleware can cue an
agent when it overlooks execution feedback~\cite{yang2026beyond}, and how an error is delivered changes
what the agent does next~\cite{woo2026same}. The choice between repairing, replanning and abandoning a
plan has been studied in theory~\cite{tewari2026ai}, repair has been compared with replanning
empirically~\cite{yuan2026replan}, and trajectory studies trace where errors enter and which controls
reach them~\cite{wang2026dependable}. Our \texttt{always} and \texttt{on\_uncertainty} arms are the
simplest such policies.

\paragraph{Replanning and self-correction.} Agents can replan from environment
feedback~\cite{sun2023adaplanner}, decompose sub-tasks they fail to execute~\cite{prasad2023adapt},
verbally reflect on task feedback~\cite{shinn2023reflexion} or refine their outputs with
self-feedback~\cite{madaan2023selfrefine}. Each of these acts on a failure signal; a silent fault
provides none. Without an external signal, models struggle to correct their own
reasoning~\cite{huang2023large}. Our triggers supply that external signal; the recovery itself is left
to the model.
